\chapter{1958年地区不明卷}

\begin{problem}
若 \(\tan \alpha=\frac{1}{2}\)，\(\tan \beta=\frac{1}{3}\)，并且 \(\alpha\) 与 \(\beta\) 均为锐角. 求证：\(\alpha+\beta=45^\circ\).
\end{problem}

\begin{solution}
由正切的和角公式，且
\(1-\tan\alpha\tan\beta=1-\frac{1}{2}\cdot\frac{1}{3}=\frac{5}{6}\ne0\)，有
\[
\tan\left(\alpha+\beta\right)
=\frac{\tan\alpha+\tan\beta}{1-\tan\alpha\tan\beta}
=\frac{\frac{1}{2}+\frac{1}{3}}{1-\frac{1}{6}}
=\frac{5}{6}\div\frac{5}{6}=1.
\]
因为 \(\alpha\)，\(\beta\) 均为锐角，所以 \(0^\circ<\alpha+\beta<180^\circ\). 在此范围内，满足 \(\tan\theta=1\) 的角只有 \(\theta=45^\circ\)，故 \(\alpha+\beta=45^\circ\).
\end{solution}

\begin{problem}
问 \(k\) 取何值时，二次方程 \(\left(2k-1\right)x^{2}+\left(k+1\right)x+k-4=0\) 才有等根.
\end{problem}

\begin{solution}
要使所给方程确为二次方程，先要求 \(2k-1\ne0\). 对二次方程来说，有等根的充要条件是判别式为零，因此
\[
0=\left(k+1\right)^2-4\left(2k-1\right)\left(k-4\right)
=k^2+2k+1-4\left(2k^2-9k+4\right)
=-7k^2+38k-15.
\]
于是 \(7k^2-38k+15=0\)，即 \((7k-3)(k-5)=0\).
所以 \(k=\frac{3}{7}\) 或 \(k=5\). 这两个值都满足 \(2k-1\ne0\)，且判别式确为零，故所求为 \(k=\frac{3}{7},5\).
\end{solution}

\begin{problem}
解方程：\(\sqrt[3]{4a^{3}+x}+\sqrt[3]{4a^{3}-x}=2a\).
\end{problem}

\begin{solution}
令 \(u=\sqrt[3]{4a^3+x}\)，\(v=\sqrt[3]{4a^3-x}\).
原方程变为 \(u+v=2a\)，并且
\[
u^3+v^3=\left(4a^3+x\right)+\left(4a^3-x\right)=8a^3.
\]
对 \(u+v=2a\) 两边立方，利用 \((u+v)^3=u^3+v^3+3uv(u+v)\)，得
\[
8a^3=8a^3+3uv\cdot2a.
\]
所以
\[
6auv=0.
\]

分情况讨论. 当 \(a=0\) 时，原方程为 \(\sqrt[3]{x}+\sqrt[3]{-x}=0\)，任意实数 \(x\) 都满足，因此此时解为任意实数 \(x\). 当 \(a\ne0\) 时，由 \(6auv=0\) 得 \(uv=0\)，所以 \(u=0\) 或 \(v=0\). 于是
\(u=0\) 时，\(4a^3+x=0\)，从而 \(x=-4a^3\)；\(v=0\) 时，\(4a^3-x=0\)，从而 \(x=4a^3\).
代回原方程检验：当 \(x=4a^3\) 时，左边为 \(\sqrt[3]{8a^3}+0=2a\)；当 \(x=-4a^3\) 时，左边为 \(0+\sqrt[3]{8a^3}=2a\). 两值均满足原方程.

故当 \(a\ne0\) 时，解为 \(x=\pm4a^3\)；当 \(a=0\) 时，任意实数 \(x\) 都是解.
\end{solution}
